\section{Resultados}
\label{results}

Esta seção apresenta a composição da base analisada, a discriminação dos
grupos na validação por vaso e a contribuição dos sensores MQ em conjunto
com as variáveis ambientais e operacionais. Todas as métricas apresentadas
referem-se às 55.946 leituras da mesma base.

\subsection{Base de dados}
\label{subsec:res_base}

A base analisada contém 55.946 leituras do nariz eletrônico, provenientes
de dez vasos: cinco inoculados com \textit{H. glycines} e cinco controles.
O grupo inoculado reúne 26.342 leituras e o controle, 29.604. Os registros
incluem os sinais dos sensores MQ, o índice de solo, a temperatura, a
pressão e a identificação do vaso e da classe experimental
(Tabela~\ref{tab:base_dados}). As leituras são observações repetidas dos
mesmos vasos e não representam 55.946 unidades biológicas independentes.

\begin{table}[H]
\centering
\footnotesize
\caption{Distribuição das leituras por grupo experimental na base analisada.}
\label{tab:base_dados}
\begin{tabular}{lrr}
\toprule
Grupo & Vasos & Leituras \\
\midrule
Inoculado (classe 0) & 5 & 26.342 \\
Controle (classe 1) & 5 & 29.604 \\
\midrule
Total & 10 & 55.946 \\
\bottomrule
\end{tabular}
\end{table}

\subsection{Discriminação dos grupos por vaso}
\label{subsec:res_loso}

Na validação \textit{leave-one-subject-out} (LOSO), cada vaso foi reservado
para avaliação em uma das dez rodadas. Com cinco sensores MQ (MQ-2, MQ-3,
MQ-8, MQ-135 e MQ-138) e as três variáveis ambientais e operacionais
(índice de solo, temperatura e pressão), o \textit{Extra Trees} (ET)
classificou corretamente 54.783 das 55.946 leituras, com acurácia de
97,92\%, acurácia balanceada de 97,88\% e F1 macro de 97,91\%
(Tabela~\ref{tab:loso_modelos}). Com apenas os cinco MQ, a acurácia do ET
foi de 94,76\% (53.015/55.946). A seleção das entradas e sua avaliação
foram feitas nos mesmos dez vasos; por isso, os resultados são
exploratórios e não constituem uma estimativa independente de desempenho
em novos vasos.

\subsection{Contribuição dos sensores MQ e das variáveis ambientais}
\label{subsec:res_mq}

Na busca de combinações de entradas do ET, o conjunto formado por índice
de solo, temperatura e pressão, sem sinais MQ, acertou 53.865 das 55.946
leituras (96,28\%). A configuração de cinco MQ com essas três variáveis
acertou 54.783 leituras (97,92\%), diferença descritiva de 1,64 ponto
percentual. Como a configuração foi escolhida e avaliada nos mesmos vasos,
essa diferença não demonstra um ganho generalizável dos MQ. O desempenho
elevado das variáveis ambientais isoladas também limita a atribuição da
separação entre os grupos exclusivamente aos sinais dos sensores MQ.

Nos seis classificadores avaliados com os mesmos cinco MQ, a inclusão
das variáveis ambientais elevou a acurácia por leitura em todos os casos
(Tabela~\ref{tab:loso_modelos}). O maior aumento ocorreu no
\textit{Random Forest}, de 84,32\% para 95,42\%; no ET, a acurácia passou
de 94,76\% para 97,92\%, e na TCN, de 88,21\% para 88,99\%. Com MQ e
variáveis ambientais, o ET apresentou a maior acurácia, seguido pela MLP
(96,07\%), pelo \textit{Random Forest} (95,42\%), pelo ensemble de cinco
modelos (95,37\%) e pelo k-NN (95,06\%). O ensemble combinou, por média
não ponderada das probabilidades, ET, \textit{Random Forest}, TCN, MLP e
k-NN; sua acurácia passou de 92,64\% com apenas MQ para 95,37\% com
MQ e variáveis ambientais. A composição do ensemble foi definida após a
inspeção dos desempenhos individuais e também deve ser tratada como
exploratória.
As diferenças entre modelos são descritivas; não foi realizado teste
pareado de superioridade entre classificadores.

\begin{table}[H]
\centering
\footnotesize
\setlength{\tabcolsep}{4pt}
\caption{Desempenho por leitura na validação LOSO por vaso com 55.946 leituras. As entradas fixas são cinco sensores MQ (MQ-2, MQ-3, MQ-8, MQ-135 e MQ-138), isolados ou acompanhados de índice de solo, temperatura e pressão. O ensemble combina Extra Trees, Random Forest, TCN, MLP e k-NN por média não ponderada das probabilidades.}
\label{tab:loso_modelos}
\begin{tabular}{lrrrr}
\toprule
Modelo & Acurácia (\%) & Acurácia bal. (\%) & Recall (\%) & F1 macro (\%) \\
\midrule
\multicolumn{5}{l}{\textit{Cinco sensores MQ}} \\
Extra Trees & 94,76 & 94,47 & 89,56 & 94,70 \\
Random Forest & 84,32 & 83,87 & 76,20 & 84,07 \\
TCN & 88,21 & 87,63 & 77,64 & 87,94 \\
MLP (ANN) & 89,58 & 89,22 & 82,95 & 89,44 \\
k-NN & 88,96 & 88,66 & 83,62 & 88,84 \\
Ensemble (5 modelos) & 92,64 & 92,21 & 84,83 & 92,52 \\
\midrule
\multicolumn{5}{l}{\textit{Cinco sensores MQ + variáveis ambientais}} \\
Extra Trees & 97,92 & 97,88 & 97,12 & 97,91 \\
Random Forest & 95,42 & 95,31 & 93,38 & 95,40 \\
TCN & 88,99 & 88,75 & 84,61 & 88,90 \\
MLP (ANN) & 96,07 & 96,08 & 96,25 & 96,06 \\
k-NN & 95,06 & 94,86 & 91,37 & 95,02 \\
Ensemble (5 modelos) & 95,37 & 95,24 & 93,04 & 95,34 \\
\bottomrule
\end{tabular}

\vspace{0.1cm}
\parbox{\linewidth}{\footnotesize Recall referente à classe 0
(inoculado).}
\end{table}
